\section{Background and Related Work}
This scoped narrative synthesis establishes context for ArchSync's design; it is not a systematic literature review. Sources were selected purposively from the retained keyword-search corpus for architecture conformance, drift detection, and static-analysis reliability, emphasizing 2022--2025. Four earlier works supply concept origins. Selection and interpretation may omit relevant approaches; coverage completeness and absence-of-prior-work claims are outside scope.

\textbf{Erosion and reconstruction.} A mapping of architecture erosion describes its manifestations and consequences \cite{li2022erosion}. Code-review studies identify symptoms in OpenStack and later OpenStack/Qt projects \cite{li2022symptoms,li2023warnings}; their shared project contexts are not independent replications. Practitioner evidence identifies documentation, responsibilities, and shared practices as ways to manage drift \cite{anthony2024drifting}. Reconstruction taxonomies organize recovery strategies \cite{ducasse2009reconstruction,sar2024taxonomy}, while a microservice mapping compares tool purposes and availability \cite{bakhtin2024mapping}. View-based analysis compares intended and reconstructed specifications in an industrial case \cite{uzun2024drift}. ArchSync restricts its observation boundary to explicit service and infrastructure relationships.

\textbf{Conformance.} Reflexion Models, static conformance comparisons, and DCL provide the origins of model comparison and allowed/forbidden dependency rules \cite{murphy1995reflexion,knodel2007comparison,terra2009dcl}; neither concept is new here. CALint applies dependency-rule checks to Python \cite{calint2022}. ArchLintor is a close TypeScript comparison, detecting module-interaction violations and supporting CI/CD \cite{archlintor2025}. ArchSync's typed service topology differs from that module boundary; no head-to-head accuracy result is claimed. IaC metrics address deployment coupling \cite{ntentos2022iac}, and 3Erefactor searches for executable inconsistency-reducing repairs \cite{refactor2024}. LLM-generated architectural tests distinguish execution success from correct rule expression \cite{lucena2025gatekeepers}. ArchSync executes explicit rules deterministically, but that choice alone implies neither accuracy nor superiority.

\textbf{Change-time feedback.} Commit-level prediction studies architectural change \cite{commits2025}; repository mining tracks microservice evolution \cite{evolution2024}; industrial conformance visualization supplies practitioner evidence \cite{thermofisher2025}. A predicted or observed change need not violate intended architecture. ArchSync instead reports findings introduced at a Git head, preserving separate violation and evolution decisions.

\textbf{Evaluation reliability.} Replication support remains an evaluation concern \cite{konersmann2022replicability}. Decomposition surveys cover an adjacent problem \cite{abgaz2023decomposition}, whereas a microservice-recovery comparison executes nine tools and reports individual and combined recovery performance \cite{schneider2025comparison}. That study supplies an external-comparison precedent, not scores directly comparable with ArchSync. Static visual models and multiple dependency models motivate future multi-source work \cite{staticmodels2024,dependencies2024}. Interrogation testing exposes both precision and soundness failures \cite{kaindlstorfer2024interrogation}, motivating our explicit hard negatives. These sources establish comparison dimensions rather than exhaustive novelty or representative accuracy.
